\booksection{Приложение к главе 1: Специальный раздел об операционной системе HarmonyOS}

\chapterintro{HarmonyOS — это операционная система, самостоятельно разработанная Huawei. С момента выпуска в 2019 году она прошла через несколько крупных версий. Это не просто «ещё одна замена Android», а совершенно новый вид системы с уникальной концепцией дизайна. Давайте вместе углубимся в изучение этой набирающей силу операционной системы.}

\booksubsection{I. Что такое HarmonyOS? Почему она особенная?}

Ключевая концепция HarmonyOS — «Интернет всего». В отличие от традиционных операционных систем, обслуживающих в основном одно устройство, HarmonyOS с самого начала проектировалась с учётом взаимодействия множества устройств.

\textbf{Распределённая архитектура}— это душа HarmonyOS. Проще говоря, она может объединять несколько устройств в виртуальный «супертерминал». Ваш телефон, планшет, компьютер, часы, телевизор — больше не изолированные устройства, а органы единого целого, которые могут использовать возможности друг друга.

Пример: вы смотрите видео на телефоне и хотите продолжить на большом экране. Традиционный способ — трансляция экрана; а в HarmonyOS вы можете напрямую «передать» видео на телевизор, а телефон станет пультом. Это не просто трансляция экрана — само приложение работает на разных устройствах.

\booksubsection{II. HarmonyOS 4.3: зрелая версия на базе AOSP}

HarmonyOS 4.3 — это версия, широко развёрнутая на текущих телефонах Huawei. Следует отметить, что HarmonyOS 4.3 совместима с кодовой базой Android Open Source Project (AOSP), поэтому она полностью совместима с приложениями Android — большинство приложений, которые вы скачиваете из магазина приложений Huawei, по сути, всё ещё являются Android-приложениями, работающими через слой совместимости HarmonyOS. Однако HarmonyOS 4.3 — это не просто оболочка поверх Android: на базе AOSP добавлено множество собственных технологий Huawei, таких как распределённые возможности, движок Ark Engine и другие.

\subsubsection*{Ключевые особенности:}

\begin{itemize}
  \item \textbf{Совместимость с Android-приложениями}: благодаря разработке на базе AOSP, HarmonyOS 4.3 может нормально запускать практически все Android-приложения, пользователи почти не замечают разницы.
  \item \textbf{Движок Ark Engine}: собственный компилятор и среда выполнения Huawei, обеспечивающие более плавную и энергоэффективную работу приложений. По сравнению с ранними версиями плавность заметно улучшилась.
  \item \textbf{Супертерминал}: в центре управления можно наглядно увидеть другие устройства HarmonyOS поблизости и подключиться перетаскиванием. Телефон может использовать клавиатуру и мышь компьютера, планшет — камеру телефона — всё очень удобно.
  \item \textbf{Сервисные карточки}: похожи на виджеты iOS, но более гибкие. Можно провести вверх по иконке приложения, чтобы развернуть карточку с ключевой информацией, не открывая приложение.
  \item \textbf{Режим чистоты}: разрешает установку приложений только из официального магазина приложений, снижая риск вредоносного ПО.
  \item \textbf{Центр приватности}: централизованное управление разрешениями приложений, просмотр записей поведения приложений.
\end{itemize}

\subsubsection*{Особенности пользовательского интерфейса:}

\begin{itemize}
  \item \textbf{Журнальный экран блокировки}: на экране блокировки отображаются красивые изображения и статьи, смахнул — и перешёл на рабочий стол.
  \item \textbf{Универсальные карточки}: иконки на рабочем столе можно развернуть в карточки с дополнительной информацией.
  \item \textbf{Супербуфер обмена}: временное хранилище, куда можно поместить текст, изображения и файлы, а затем перетащить в любое приложение. Межприложенческий обмен становится невероятно плавным.
  \item \textbf{Умный поиск}: глобальный поиск — можно искать приложения, файлы, настройки и веб-страницы.
\end{itemize}

\booksubsection{III. HarmonyOS NEXT (HarmonyOS 6): настоящая точка отсчёта собственной ОС Huawei}

HarmonyOS NEXT (HarmonyOS NEXT 5.0) имеет фундаментальное отличие от HarmonyOS 4.3 — она использует собственное микроядро HarmonyOS, больше не зависит от ядра Linux и кода AOSP, а также\textbf{больше не поддерживает напрямую Android-приложения}. Это ключевой шаг перехода HarmonyOS от «переходной формы, совместимой с Android» к «полностью независимой операционной системе».

\subsubsection*{Важные изменения:}

\begin{itemize}
  \item \textbf{Полностью собственное ядро}: HarmonyOS NEXT использует самостоятельно разработанное Huawei ядро и больше не содержит никакого кода AOSP. Система стала компактнее, безопаснее и эффективнее.
  \item \textbf{Больше не поддерживает Android-приложения}: это самое большое изменение. HarmonyOS NEXT может запускать только нативные HarmonyOS-приложения и не может устанавливать и запускать APK-файлы. Перед обновлением убедитесь, что у ваших часто используемых приложений уже есть нативные версии для HarmonyOS.
  \item \textbf{Экосистема нативных приложений}: Huawei активно продвигает разработчиков к переносу приложений на нативную платформу HarmonyOS. В настоящее время основные приложения, такие как WeChat, Alipay, Taobao и другие, уже выпустили нативные версии для HarmonyOS, но некоторые нишевые приложения могут ещё не быть адаптированы.
  \item \textbf{Архитектура безопасности Star Shield}: совершенно новая система безопасности с полной цепочкой защиты от аппаратного до программного уровня. Поддерживает более детальный контроль разрешений, например «размытое местоположение» — приложение знает только, где вы примерно находитесь, но не точное местоположение.
  \item \textbf{Глубокая интеграция ИИ}: ИИ-возможности системного уровня, а не только в отдельных приложениях. Например, ИИ-субтитры могут переводить видео в реальном времени, ИИ-выделение объектов может умно распознавать субъекты.
  \item \textbf{Бесшовная передача между устройствами}: дальнейшее развитие распределённых возможностей. Состояние приложения может бесшовно переключаться между устройствами — документ, который вы редактируете на телефоне, автоматически открывается на компьютере, когда вы подходите к нему.
  \item \textbf{Мета-сервисы}: лёгкие сервисы, не требующие установки. Например, сканирование QR-кода напрямую открывает определённую функцию без необходимости сначала скачивать приложение.
\end{itemize}

\begin{tcbwarn}{Обязательно прочтите перед обновлением до HarmonyOS NEXT}
После обновления до HarmonyOS NEXT ваш телефон не сможет запускать Android-приложения. Обязательно перед обновлением убедитесь: (1) у ваших часто используемых приложений уже есть нативные версии HarmonyOS; (2) в вашей работе или жизни нет сценариев, обязательно требующих Android-приложений. После обновления откат до HarmonyOS 4.3 невозможен.

\end{tcbwarn}

\subsubsection*{Изменения в логике управления:}

\begin{itemize}
  \item \textbf{Жестовая навигация}: смахивание от левого/правого края — назад, смахивание снизу вверх — на рабочий стол, смахивание вверх с удержанием — многозадачность.
  \item \textbf{Умная многозадачность}: поддержка разделённого экрана, плавающих окон, параллельного просмотра и других форм многозадачности.
  \item \textbf{Звонки ChangLian}: сервис видеозвонков Huawei с поддержкой демонстрации экрана, улучшения лица и размытия фона.
\end{itemize}

\booksubsection{IV. Экосистема HarmonyOS: устройства и приложения}

HarmonyOS — это не только телефонная система, а экосистема, охватывающая множество типов устройств.

\subsubsection*{Поддерживаемые типы устройств:}

\begin{itemize}
  \item \textbf{Телефоны и планшеты}: основные устройства, наиболее полная функциональность.
  \item \textbf{Умные часы и браслеты}: работают независимо или совместно с телефоном.
  \item \textbf{Умные экраны (телевизоры)}: центр домашних развлечений.
  \item \textbf{Умные колонки}: точка входа для голосового взаимодействия.
  \item \textbf{Автомобили}: кабина HarmonyOS, опыт интеллектуального вождения.
  \item \textbf{Умный дом}: холодильники, кондиционеры, стиральные машины и т.д.
\end{itemize}

\subsubsection*{Текущее состояние экосистемы приложений:}

В магазине приложений HarmonyOS уже миллионы приложений, основные из них в основном присутствуют. Huawei также активно продвигает разработчиков к созданию нативных приложений HarmonyOS. По состоянию на середину 2026 года такие ведущие приложения, как WeChat, Alipay, Taobao, Douyin, Meituan, JD.com, уже выпустили нативные версии для HarmonyOS, повседневное использование практически безбарьерно. Некоторые нишевые и отраслевые приложения могут всё ещё находиться в процессе адаптации; перед обновлением рекомендуется проверить, поддерживаются ли ваши часто используемые приложения.

\booksubsection{V. Практические советы по HarmonyOS}

\begin{itemize}
  \item \textbf{Супертерминал}: смахните вниз из правого верхнего угла экрана, чтобы открыть центр управления, нажмите «Супертерминал», чтобы увидеть другие устройства поблизости, и подключитесь перетаскиванием.
  \item \textbf{Супербуфер обмена}: долгое нажатие на текст/изображение/файл, выберите «Копировать в буфер», затем в другом приложении долгое нажатие для вставки или перетащите оттуда.
  \item \textbf{Сервисные карточки}: смахните вверх по иконке приложения, чтобы развернуть карточку; долгое нажатие на карточку позволяет изменить размер или добавить на рабочий стол.
  \item \textbf{Умная многозадачность}: смахните от левого или правого края экрана внутрь с удержанием, чтобы открыть боковую панель; нажмите на приложение, чтобы открыть его в плавающем окне.
  \item \textbf{Защита приватности}: Настройки > Приватность, где можно просмотреть записи использования разрешений приложений, настроить размытое местоположение и т.д.
\end{itemize}

\begin{tcbwarn}{Примечания}
HarmonyOS постоянно обновляется и итеративно развивается, разные версии могут иметь некоторые отличия. Описание в этой статье основано на основных характеристиках HarmonyOS 4.3 и 6; конкретные функции зависят от фактической версии вашего телефона. Рекомендуется поддерживать систему в актуальном состоянии для получения новейших функций и исправлений безопасности.

\end{tcbwarn}

\begin{tcbpractice}{Попробуйте прямо сейчас: испытайте уникальные функции HarmonyOS}
\begin{enumerate}
  \item Откройте центр управления, найдите «Супертерминал» и посмотрите, есть ли поблизости другие устройства HarmonyOS.
  \item Попробуйте «Супербуфер обмена»: скопируйте текст в буфер, затем перетащите его в заметки.
  \item Смахните вверх по иконкам нескольких часто используемых приложений и посмотрите, какое содержание отображается в их сервисных карточках.
  \item Перейдите в Настройки > Приватность и посмотрите, какие приложения недавно использовали чувствительные разрешения.
\end{enumerate}

\end{tcbpractice}
